\subsection{Refinement on the Physical Robot}
\label{sec:exp-real-robot}

The simulation comparisons isolate the actor update; the physical study
asks whether it adds value when refining an existing robot skill through
practice. Success-only BC provides a critic-free alternative, separating
the benefit of imitating successful attempts from that of learning from
both outcomes. We use an OpenArm platform's left
seven-joint arm and gripper, with chest and left-wrist RGB cameras, measured
joint and gripper positions, and a fixed task-language prompt. The two
tasks are \emph{Jigglypuff into the case} and \emph{red-on-purple stacking}.
Both use a hand-made staircase: different support heights and nearby edges
make approach, grasp alignment, and release important. Stacking additionally
requires accurate placement before withdrawing the gripper.

\textbf{Four-way comparison.}
For each task, a flow-matching policy is pretrained on teleoperated
demonstrations. It predicts 30-step chunks of absolute joint and gripper
targets at 15 Hz using a ten-step Euler sampler; deployment executes a
prefix before replanning. All refinement conditions start from the same
checkpoint. Success-only BC continues supervised training on successful
online trajectories. DICE-RL and DICE-RL + CAST use both successful and
failed attempts, with the same frozen base, critic architecture, replay,
exploration, and update schedule. This tests whether critic-guided
refinement adds value beyond imitating successful practice, and whether
the actor-update rule changes that benefit.

\textbf{Collection and evaluation.}
After each attempt, an operator assigns the success/failure label. Valid
failures are retained for the RL arms, and invalid or discarded trials are
accounted for explicitly. Replay records commands actually executed after
chunk-prefix selection and command limiting. The refinement protocol uses
matched robot interaction budgets and execution limits; evaluation attempts
are held out from refinement data. Each method is evaluated over 20 rollouts
per task. For Jigglypuff, the toy's pose is randomly changed on the stairs
between attempts. Exact pose ranges and paired resets across methods are
not established by the available evaluation record. Toy placement
requires release fully inside the case with the toy remaining there.
Stacking requires the red cube to remain supported by the purple cube
after release and gripper withdrawal; the post-release observation interval
is not specified in the reported stacking results.

\textbf{Results across both tasks.}
Fig.~\ref{fig:real-robot} reports 11/20 successes for the base FM policy,
16/20 after success-only BC refinement, 18/20 for DICE-RL, and 19/20 for
DICE-RL + CAST on Jigglypuff. On stacking, the same methods achieve
7/20, 13/20, 12/20, and 14/20 successes. All refinements improve observed
success over the base on both tasks. CAST reaches the highest observed
rates, 95\% and 70\%, gaining eight and seven successes over the bases.
The comparison with BC is task-dependent: DICE-RL exceeds BC on Jigglypuff
but trails it on stacking. CAST exceeds DICE-RL by one and two trials,
and BC by three and one. These small margins support further evaluation
rather than a reliable ranking of refinement methods.
Fig.~\ref{fig:real-robot} illustrates both tasks using separate
demonstration sequences.
